
\begin{tabularx}{\textwidth}{L{30mm}L{24mm}L{32mm}X}
\toprule
\textbf{Változat} & \textbf{Sebesség} & \textbf{Közeg} & \textbf{Jellemző} \\
\midrule
10BASE5 & 10 Mbit/s & vastag koax & Klasszikus sárga kábel, busz jelleg, kb. 500 m szegmenshossz. \\
10BASE2 & 10 Mbit/s & vékony koax & Cheapernet, BNC és T-csatlakozók, kb. 185 m szegmenshossz. \\
10BASE-T & 10 Mbit/s & csavart érpár & Pont--pont kábelezés hubhoz vagy switchhez, jellemzően 100 m. \\
100BASE-TX/FX & 100 Mbit/s & UTP vagy optika & Fast Ethernet, 802.3u, kapcsolt hálózatokban nagyon elterjedt. \\
1000BASE-SX/LX/T & 1 Gbit/s & optika vagy UTP & Gigabit Ethernet; 1000BASE-T tipikusan 100 m Cat5 vagy jobb kábelen. \\
10GBASE-* & 10 Gbit/s & főleg optika, később réz & 10 Gigabit Ethernet; full duplex működés, a keretformátum alapvetően változatlan. \\
\bottomrule
\end{tabularx}
